\chapter{INTRODUÇÃO}
\label{chap:contextualizacao}

Todo TCC começa com uma Introdução, aquela seção onde você tenta convencer o leitor de que o seu tema é relevante e indispensável para a humanidade - mesmo que seja só um sistema que manda notificação quando a pizza sai do forno.

A introdução também serve pode contar uma história emocionante: do problema que ninguém resolvia, até a sua solução (que surgiu às 3h da manhã, enquanto você rolava na cama tentando dormir).

Você descreve o contexto, apresenta o desafio, dramatiza um pouco e termina com uma pergunta do tipo: \enquote{Diante disso, como garantir a eficiência da metodologia proposta em cenários reais?} Você não faz ideia. O seu trabalho é justamente encontrar a resposta.

Exemplo de como incluir uma ilustração: primeiro, explique e dê o contexto da ilustração, e a referencie, gráfico \ref{graph:anos_trabalhos_selecionados_v2}.

\input{images/resultados/anos_trabalhos_selecionados_v2}

Bla bla bla bla figura \ref{fig:figura_panda}.

\input{images/resultados/figura_panda}

Por fim, não se esqueça de incluir a pergunta norteadora: \textit{\enquote{Como otimizar o processo de aquecimento de pizza utilizando softwares com sistemas de notificação?}} Ou melhor: Será que se eu juntar esses negócio aqui, vai dar bom?

Bla bla bla bla diagrama \ref{diag:fluxo_addon}

\input{images/resultados/fluxo_addon}

% \begin{grafico}[ht!]
% \centering

% \caption{\textmd{Texto Texto Texto Texto Texto Texto Texto}}
% \label{graph:grafico1}
% \fcolorbox{gray}{white}{\includegraphics[width=0.80\textwidth]{images/captitulo1/figuraex.png}}

% \fonteref{\citeauthor{akhtar2018threat} (\citeyear{akhtar2018threat})}
% \end{grafico}

% Texto texto texto texto texto texto texto texto texto texto texto texto texto texto texto texto texto texto texto texto texto texto texto texto texto texto texto \gls{UFPE} ver Diagrama \ref{diag:diagrama1}.

% \begin{diagrama}[ht!]
% \centering

% \caption{\textmd{Texto do diagrama}}
% \label{diag:diagrama1}
% \fcolorbox{gray}{white}{\includegraphics[width=0.80\textwidth]{images/captitulo1/figuraex2.png}}

% \fonteadapt{\citeauthor{akhtar2018threat} (\citeyear{akhtar2018threat})}
% \end{diagrama}
    